\documentclass[11pt]{article}

\usepackage[margin=1in]{geometry}
\usepackage{amsmath,amssymb,amsthm,mathrsfs,graphicx}
\usepackage{tikz-cd}
\usepackage{quiver}
\usepackage[colorlinks=true,linkcolor=blue,urlcolor=blue,citecolor=blue]{hyperref}

\providecommand{\A}{}
\renewcommand{\A}{\mathbb{A}}
\providecommand{\C}{}
\renewcommand{\C}{\mathbb{C}}
\providecommand{\F}{}
\renewcommand{\F}{\mathcal{F}}
\providecommand{\G}{}
\renewcommand{\G}{\mathcal{G}}
\providecommand{\bb}{}
\renewcommand{\bb}{\mathbb}
\providecommand{\td}{}
\renewcommand{\td}{\tilde}
\providecommand{\epsi}{}
\renewcommand{\epsi}{\varepsilon}
\providecommand{\exist}{\exists}
\providecommand{\mf}{}
\renewcommand{\mf}{\mathfrak}
\providecommand{\bs}{}
\renewcommand{\bs}{\backslash}
\providecommand{\supp}{}
\renewcommand{\supp}{\operatorname{Supp}}
\providecommand{\del}{}
\renewcommand{\del}{\partial}
\providecommand{\wg}{}
\renewcommand{\wg}{\wedge}
\providecommand{\proj}{}
\renewcommand{\proj}{\operatorname{proj}}
\providecommand{\mc}[1]{}
\renewcommand{\mc}[1]{\mathcal{#1}}
\providecommand{\ms}[1]{}
\renewcommand{\ms}[1]{\mathscr{#1}}
\providecommand{\O}{}
\renewcommand{\O}{\mathcal{O}}
\providecommand{\op}[1]{}
\renewcommand{\op}[1]{\operatorname{#1}}
\providecommand{\P}{}
\renewcommand{\P}{\mathbb{P}}
\providecommand{\Q}{}
\renewcommand{\Q}{\mathbb{Q}}
\providecommand{\R}{}
\renewcommand{\R}{\mathbb{R}}
\providecommand{\Z}{}
\renewcommand{\Z}{\mathbb{Z}}
\DeclareMathOperator{\Spec}{Spec}
\DeclareMathOperator{\Hom}{Hom}
\providecommand{\id}{}
\renewcommand{\id}{\operatorname{id}}
\providecommand{\im}{}
\renewcommand{\im}{\operatorname{im}}
\providecommand{\sgn}{}
\renewcommand{\sgn}{\operatorname{sgn}}
\providecommand{\rad}{}
\renewcommand{\rad}{\operatorname{rad}}
\providecommand{\md}{}
\renewcommand{\md}{\operatorname{mod}}
\providecommand{\sign}{}
\renewcommand{\sign}{\operatorname{sign}}
\providecommand{\Ann}{}
\renewcommand{\Ann}{\operatorname{Ann}}
\providecommand{\End}{}
\renewcommand{\End}{\operatorname{End}}
\providecommand{\Aut}{}
\renewcommand{\Aut}{\operatorname{Aut}}
\providecommand{\coker}{}
\renewcommand{\coker}{\operatorname{coker}}
\providecommand{\spa}{}
\renewcommand{\spa}{\operatorname{span}}
\providecommand{\gr}{}
\renewcommand{\gr}{\operatorname{Gr}}
\providecommand{\ot}{}
\renewcommand{\ot}{\otimes}
\providecommand{\range}{}
\renewcommand{\range}{\text{Range}}
\providecommand{\nb}{}
\renewcommand{\nb}{\nabla}
\providecommand{\bk}[1]{}
\renewcommand{\bk}[1]{\langle #1\rangle}
\providecommand{\bwg}{}
\renewcommand{\bwg}{\bigwedge\nolimits}
\providecommand{\Tor}{}
\renewcommand{\Tor}{\operatorname{Tor}}
\providecommand{\Bl}{}
\renewcommand{\Bl}{\operatorname{Bl}}

\title{Solutions to Vakil's The Rising Sea}
\author{Tianyi Wang}
\date{\today}

\begin{document}
\maketitle
\tableofcontents

\section{1 -- 1.1.B -- Automorphism Group}
% id: 1-1-1-b-automorphism-group
% author: Tianyi
% updated: 2026-08-27

Let $A$ be an object in a category $\ms C$. The invertible elements $\op{Aut}(A)\subset \op{Mor}(A,A)$ indeed forms a group because we have the identity $\op{id}_A: A\to A$, and every morphism is invertible by definition. Moreover associativity holds by definition of a category. Hence $\op{Aut}(A)$ is indeed a group.

In the example of sets, $\op{Aut}(A)$ is the set of bijections from $A$ to itself; In the example of vector spaces, $\op{Aut}(V)$ is the set of linear isomorphisms from $V$ to itself.

\section{1 -- 1.1.C -- Double Dual Functor}
% id: 1-1-1-c-double-dual-functor
% author:
% updated: 2026-08-27

This amounts to checking commutativity of the following diagram:
% https://q.uiver.app/#q=WzAsNCxbMCwwLCJWXntcXHZlZVxcdmVlfSJdLFsxLDAsIldee1xcdmVlXFx2ZWV9Il0sWzAsMSwiViJdLFsxLDEsIlciXSxbMCwxLCJmXntcXHZlZVxcdmVlfSJdLFsyLDMsImYiLDJdLFsyLDAsIm1fViJdLFszLDEsIm1fVyIsMl1d
\[\begin{tikzcd}[cramped]
	{V^{\vee\vee}} & {W^{\vee\vee}} \\
	V & W
	\arrow["{f^{\vee\vee}}", from=1-1, to=1-2]
	\arrow["{m_V}", from=2-1, to=1-1]
	\arrow["f"', from=2-1, to=2-2]
	\arrow["{m_W}"', from=2-2, to=1-2]
\end{tikzcd}\]
Here, the map $m_V: V^{\vee\vee}\to V$ is defined to be
$$
m_V:v\mapsto \epsilon_v
$$
where $\epsilon_v: V^\vee\to k$ is defined to be the delta mass at the funcional $\delta_v\in V^\vee$, where $\delta_v: V\to k$ is the delta mass at the vector $v\in V$.

\section{1 -- 1.1.D -- Dimensions and Basis Characterize Vector Spaces}
% id: 1-1-1-d-dimensions-characterizes-vector-spaces
% author: Tianyi
% updated: 2026-08-27

We show that $\ms V\to \op{f.d.Vec}_k$ is an equivalence of categories. We construct the inverse functor as follows: Simultaneously choose basis for each vector space $V\in \op{f.d.Vec}_k$ (which Vakil says we can do). Define the functor $\rho: \op{f.d.Vec}_k\to \ms V$ as follows:

Let $V\in \op{f.d.Vec}_k$ be a vector space with chosen basis $e_1,...,e_n$. Then define $\rho(V)=k^n$. If $W$ is another vector space with basis $f_1,...,f_m$, so $\rho(W)=k^m$, and if $T:V\to W$ is a linear map, then with respect to these basis $T$ can be identified with an $m\times n$ matrix, which is a linear map $\rho(T):k^n\to k^m$.

\section{1 -- 1.2.A -- Initial and Final Objects}
% id: 1-1-2-a-initial-and-final-objects
% author: Tianyi
% updated: 2026-08-28

Let $\ms C$ be a category and $I,I'$ be two initial objects in the category. Then there is a unique map $f: I\to I'$ and a unique map $g: I'\to I$. So we get a map $g\circ f:I\to I$. But we also have $\op{id}_I: I\to I$. Hence we conclude that $g\circ f=\op{id}_I$ and similarly $f\circ g=\op{id}_{I'}$. Hence $I,I'$ are isomorphic. Similar argument shows that two final objects are isomorphic.

\section{1 -- 1.2.B -- Untitled solution}
% id: 1-1-2-b
% author: Tianyi
% updated: 2026-08-28

In both $Sets$ and $Tops$, the initial object is $\emptyset$.

There is precisely one map $\emptyset\to X$ for any object $X$, Since there are no elements in $\emptyset$ whose images need to be specified.

The final object is a singleton $\{*\}$.

There is precisely one map $X\to\{*\}$ which is the constant map.

In $Rings$, both the initial and final object is the 0 ring.

\section{1 -- 1.2.C -- $A\to S^{-1}A$ injective iff no zero divisors in S}
% id: 1-1-2-c-a-to-s-1-a-injective-iff-no-zero-divisor
% author: Tianyi
% updated: 2026-08-28

Suppose $S$ has a zero divisor $a$. Then there exists a nonzero $b$ such that $ab=0$. Then $b/1=0/1$ because $a(b\cdot 1-0\cdot 1)=0$. But $b$ is nonzero, so the map $A \to S^{-1}A$ is not injective.

Conversely, suppose the map $A\to S^{-1}A$ is not injective, which means there is a nonzero $a$ such that $a/1=0/1$. Then there exists $s\in S$ such that $s(a-0)=0$. Hence $sa=0$ ans $s$ is a zero divisor.

\section{1 -- 1.2.D -- Universal Property of Localization}
% id: 1-1-2-d-universal-property-of-localization
% author: Tianyi
% updated: 2026-08-28

We want to show that any map $\phi: A \to B$ that sends $S$ to invertible elements in $B$ factors through the localization $A\to S^{-1}A$, i.e., the following diagram commutes:
% https://q.uiver.app/#q=WzAsMyxbMCwwLCJBIl0sWzIsMCwiU157LTF9QSJdLFsxLDEsIkIiXSxbMCwyLCJcXHBoaSIsMl0sWzAsMV0sWzEsMiwiXFxleGlzdHMhXFxwc2kiLDAseyJzdHlsZSI6eyJib2R5Ijp7Im5hbWUiOiJkYXNoZWQifX19XV0=
\[\begin{tikzcd}[cramped]
	A && {S^{-1}A} \\
	& B
	\arrow[from=1-1, to=1-3]
	\arrow["\phi"', from=1-1, to=2-2]
	\arrow["{\exists!\psi}", dashed, from=1-3, to=2-2]
\end{tikzcd}\]
Indeed, define $\psi: S^{-1}A\to B$ via $\psi(a/s)=\phi(a)\phi(s)^{-1}$. It is easy to check that this construction makes the diagram commute and is unique.

\section{1 -- 1.2.E -- Localization with Sums and Products}
% id: 1-1-2-e-localization-with-sums-and-products
% author: Tianyi
% updated: 2026-08-29

For (a) and (b) we prove the isomorphism exists by the universal property of localization. That is, given $\phi$ sending $S$ to invertible elements of $N$ we want to show
% https://q.uiver.app/#q=WzAsMyxbMCwwLCJNXzFcXHRpbWVzXFxjZG90c1xcdGltZXMgTV9uIl0sWzIsMCwiU157LTF9TV8xXFx0aW1lc1xcY2RvdHNcXHRpbWVzIFNeey0xfU1fbiJdLFsxLDEsIk4iXSxbMCwyLCJcXHBoaSIsMl0sWzAsMV0sWzEsMiwiXFxleGlzdHMhXFxwc2kiLDAseyJzdHlsZSI6eyJib2R5Ijp7Im5hbWUiOiJkYXNoZWQifX19XV0=
\[
\begin{tikzcd}[cramped]
	{M_1\times\cdots\times M_n} &  & {S^{-1}M_1\times\cdots\times S^{-1}M_n} \\
	 & {N} &
	\arrow["{\phi}"', from=1-1, to=2-2]
	\arrow[from=1-1, to=1-3]
	\arrow["{\exists!\psi}", dashed, from=1-3, to=2-2]
\end{tikzcd}
\]
where the verticle map is the obvious one $(m_1,...,m_n)\mapsto (m_1/1,...,m_n/1).$ Now let $e_1,...,e_n$ denote the standard basis in $M_1\times\cdots\times M_n$. Then we see that we can define
\begin{align*}
\psi(m_1/s_1,...,m_n/s_n)&=\psi(\sum_{1\leq i\leq n}(0,...,m_i/s_i,...,0))=\sum_{1\leq i\leq n}\phi(s_i)^{-1}\phi(0,...,m_i,...,0).
\end{align*}

(b) Recall that $\bigoplus_{i\in I}M_i$ is the collection of elements of the form $(m_i)_{i\in I}$ where $m_i=0$ for all but finitely many $i$. Again similar argument shows that we have a commutative diagram
% https://q.uiver.app/#q=WzAsMyxbMCwwLCJcXGJpZ29wbHVzX3tpXFxpbiBJfSBNX2kiXSxbMiwwLCJcXGJpZ29wbHVzX3tpXFxpbiBJfSBTXnstMX1NX2kiXSxbMSwxLCJOIl0sWzAsMiwiXFxwaGkiLDJdLFswLDFdLFsxLDIsIlxcZXhpc3RzIVxccHNpIiwwLHsic3R5bGUiOnsiYm9keSI6eyJuYW1lIjoiZGFzaGVkIn19fV1d
\[
\begin{tikzcd}[cramped]
	{\bigoplus_{i\in I} M_i} &  & {\bigoplus_{i\in I} S^{-1}M_i} \\
	 & {N} &
	\arrow["{\phi}"', from=1-1, to=2-2]
	\arrow[from=1-1, to=1-3]
	\arrow["{\exists!\psi}", dashed, from=1-3, to=2-2]
\end{tikzcd}
\]
and the map $\psi$ is defined similarly as in (a). Note that the reason that we can define it using a finite sum is because $m_i=0$ for all but finitely many $i$.

(c) Take the index set $I=\mathbb N$ and $M_i=\Z$. Let $S=\Z-\{0\}$. Then $\prod_{i\in I}S^{-1}M_i=\prod_{i\in I}\Q$. But note that $(1,1/2,1/3,1/4,...)\in \prod_{i\in I}\Q$ but this element is not in $S^{-1}\prod_{i\in I}M_i=S^{-1}\prod_{i\in I}\Z$ because we cannot find the common denominator of the said element $(1,1/2,...)$ in $S$ since it is infinite.

\section{1 -- 1.2.G -- $\Z/(10)\otimes_{\Z}\Z/(12)\cong \Z/(2).$}
% id: 1-1-2-g-z-10-otimes-z-z-12-cong-z-2
% author: Tianyi
% updated: 2026-08-29

This is actually an interesting problem. Note that since tensor product is bilinear, everything is actually generated by the element $x=\bar 1\otimes \bar 1$, since $\bar a\otimes \bar b=(ab)x$. Now since $10x=\overline {10}\otimes \bar 1=0$ and $12x=\bar 1\otimes \overline{12}=0$, we have
\[
0=12x-10x=2x.
\]
Hence $x$ has order divisible by 2. Therefore, the ring $\Z/(10)\otimes_{\Z}\Z/(12)$ must be isomorphic to $\Z/(2)$.

\section{1 -- 1.2.H -- Tensor Product is Right Exact}
% id: 1-1-2-h-tensor-product-is-right-exact
% author: Tianyi
% updated: 2026-08-29

Suppose we have an exact sequence of $A$-modules
\[
M'\xrightarrow{f} M\xrightarrow{g} M''\to 0
\]
Tensor product functor $\otimes _A N$, denoted by $\otimes N$ for convienience gives a sequence
\[
M'\ot N\xrightarrow{f\ot 1} M\ot N\xrightarrow{g\ot 1} M''\ot N\to 0.
\]
We use the shorthand $F:=f\ot 1, G:=g\ot 1$.

Since $\im f=\ker g$, it follows immediately that $\im F\subset \ker G$.

Also, since $g$ is surjective, so is $G$: Suppose we have a pure tensor $m''\ot n\in M''\ot N$, then by surjectivity there is an $m\in M$ such that $g(m)=m''$. Hence $G(m\ot n)=m''\ot n$, and surjectivity of $G$ follows from linearity.

Hence we only need to prove $\ker G\subset \im F$. Note that we have an map
\[G^\#:M\ot N/\im F\to M''\ot N\]
induced by $G$. We claim that this is actually an isomorphism by finding an inverse. First, We construct a map
\[
M''\times N\to M\ot N/\im F
\]
as follows: For every $(m'',n)\in M''\times N$, surjectivity of $g$ gives an element $m\in M$ such that $g(m)=m''$. Then the map is $(m'',n)\mapsto m\ot n\op{mod }\im F$. Then it follows from universal property of tensor product that we have a map
\[
M''\ot N\to M\ot N/\im F
\]
which one can check is an inverse to $G^{\#}$. Since $G^{\#}$ is an isomorphism and is induced from $G$, it follows that $\ker G=\im F$, as desired.

\section{1 -- 1.2N -- Fibered Products of Sets}
% id: 1-1-2n-fibered-products-of-sets
% author: Tianyi
% updated: 2026-08-29

Let $X\times_Z Y=\{(x,y)\in X\times Y: \alpha(x)=\beta(y)\}$. Suppose we have the following diagram
% https://q.uiver.app/#q=WzAsNSxbMSwxLCJYXFx0aW1lc19aIFkiXSxbMSwyLCJYIl0sWzIsMSwiWSJdLFsyLDIsIloiXSxbMCwwLCJXIl0sWzEsMywiXFxhbHBoYSJdLFsyLDMsIlxcYmV0YSIsMl0sWzAsMV0sWzAsMl0sWzQsMSwiXFxwaGkiLDIseyJjdXJ2ZSI6MX1dLFs0LDIsIlxccHNpIiwwLHsiY3VydmUiOi0yfV0sWzQsMCwiXFxyaG8iLDAseyJzdHlsZSI6eyJib2R5Ijp7Im5hbWUiOiJkYXNoZWQifX19XV0=
\[
\begin{tikzcd}[cramped]
	{W} &  &  \\
	 & {X\times_Z Y} & {Y} \\
	 & {X} & {Z}
	\arrow["{\alpha}", from=3-2, to=3-3]
	\arrow["{\beta}"', from=2-3, to=3-3]
	\arrow[from=2-2, to=3-2]
	\arrow[from=2-2, to=2-3]
	\arrow["{\phi}"', curve={height=6pt}, from=1-1, to=3-2]
	\arrow["{\psi}", curve={height=-12pt}, from=1-1, to=2-3]
	\arrow["{\rho}", dashed, from=1-1, to=2-2]
\end{tikzcd}
\]
We can define the map $\rho: W\to X\times_Z Y$ as $\rho(w)=(\phi(w),\psi(w)).$ It is straightforward to check that the map is well-defined (i.e., lands in the fiber product), and is unique.

\section{1 -- 1.2O -- Fibered Products in the Category of Open Sets}
% id: 1-1-2o-fibered-products-in-the-category-of-open-
% author: Tianyi
% updated: 2026-08-29

Fiber Product is simply intersection of open sets in this category.

\section{1 -- 1.2.P -- Terminal Fiber Product is Product}
% id: 1-1-2-p-terminal-fiber-product-is-product
% author: Tianyi
% updated: 2026-08-29

Another beautiful Saturday wasted on universal property... Let $Z$ be the terminal object in $\ms C$. Let $\mu:X\times Y\to X$ and $\nu: X\times Y\to Y$ be the projection maps of the product. Since $Z$ is terminal, universal property gives the following two commutative diagrams with maps $\phi,\psi$:

% https://q.uiver.app/#q=WzAsMTAsWzEsMSwiWFxcdGltZXMgWSJdLFsxLDIsIlgiXSxbMiwxLCJZIl0sWzIsMiwiWiJdLFswLDAsIlhcXHRpbWVzX1pZIl0sWzUsMSwiWFxcdGltZXNfWiBZIl0sWzYsMSwiWSJdLFs1LDIsIlgiXSxbNiwyLCJaIl0sWzQsMCwiWFxcdGltZXMgWSJdLFswLDIsIlxcbnUiXSxbMCwxLCJcXG11IiwyXSxbMSwzXSxbMiwzXSxbNCwxLCJwcl9YIiwyLHsiY3VydmUiOjJ9XSxbNCwyLCJwcl9ZIiwwLHsiY3VydmUiOi0yfV0sWzQsMCwiXFxleGlzdHMhXFxwaGkiLDAseyJzdHlsZSI6eyJib2R5Ijp7Im5hbWUiOiJkYXNoZWQifX19XSxbNyw4XSxbNiw4XSxbNSw3LCJwcl9YIiwyXSxbNSw2LCJwcl9ZIl0sWzksNywiXFxtdSIsMix7ImN1cnZlIjoyfV0sWzksNiwiXFxudSIsMCx7ImN1cnZlIjotMn1dLFs5LDUsIlxcZXhpc3QhXFxwc2kiLDAseyJzdHlsZSI6eyJib2R5Ijp7Im5hbWUiOiJkYXNoZWQifX19XV0=
\[
\begin{tikzcd}[cramped]
	{X\times_ZY} &  &  &  & {X\times Y} &  &  \\
	 & {X\times Y} & {Y} &  &  & {X\times_Z Y} & {Y} \\
	 & {X} & {Z} &  &  & {X} & {Z}
	\arrow["{\nu}", from=2-2, to=2-3]
	\arrow["{\mu}"', from=2-2, to=3-2]
	\arrow[from=3-2, to=3-3]
	\arrow[from=2-3, to=3-3]
	\arrow["{pr_X}"', curve={height=12pt}, from=1-1, to=3-2]
	\arrow["{pr_Y}", curve={height=-12pt}, from=1-1, to=2-3]
	\arrow["{\exists!\phi}", dashed, from=1-1, to=2-2]
	\arrow[from=3-6, to=3-7]
	\arrow[from=2-7, to=3-7]
	\arrow["{pr_X}"', from=2-6, to=3-6]
	\arrow["{pr_Y}", from=2-6, to=2-7]
	\arrow["{\mu}"', curve={height=12pt}, from=1-5, to=3-6]
	\arrow["{\nu}", curve={height=-12pt}, from=1-5, to=2-7]
	\arrow["{\exist!\psi}", dashed, from=1-5, to=2-6]
\end{tikzcd}
\]

The claim is that $\phi,\psi$ are inverses of each other hence establishes isomorphisms $X\times Y\cong X\times_Z Y$. To see this, compose the two diagrams:
% https://q.uiver.app/#q=WzAsNixbMiwyLCJYXFx0aW1lcyBZIl0sWzIsMywiWCJdLFszLDIsIlkiXSxbMywzLCJaIl0sWzEsMSwiWFxcdGltZXNfWlkiXSxbMCwwLCJYXFx0aW1lcyBZIl0sWzAsMiwiXFxudSJdLFswLDEsIlxcbXUiLDJdLFsxLDNdLFsyLDNdLFs0LDEsInByX1giLDEseyJjdXJ2ZSI6Mn1dLFs0LDIsInByX1kiLDEseyJjdXJ2ZSI6LTJ9XSxbNCwwLCJcXHBoaSJdLFs1LDEsIlxcbXUiLDAseyJjdXJ2ZSI6NX1dLFs1LDIsIlxcbnUiLDIseyJjdXJ2ZSI6LTV9XSxbNSw0LCJcXHBzaSJdXQ==
\[
\begin{tikzcd}[cramped]
	{X\times Y} &  &  &  \\
	 & {X\times_ZY} &  &  \\
	 &  & {X\times Y} & {Y} \\
	 &  & {X} & {Z}
	\arrow["{\nu}", from=3-3, to=3-4]
	\arrow["{\mu}"', from=3-3, to=4-3]
	\arrow[from=4-3, to=4-4]
	\arrow[from=3-4, to=4-4]
	\arrow["{pr_X}", curve={height=12pt}, from=2-2, to=4-3]
	\arrow["{pr_Y}", curve={height=-12pt}, from=2-2, to=3-4]
	\arrow["{\phi}", from=2-2, to=3-3]
	\arrow["{\mu}", curve={height=30pt}, from=1-1, to=4-3]
	\arrow["{\nu}"', curve={height=-30pt}, from=1-1, to=3-4]
	\arrow["{\psi}", from=1-1, to=2-2]
\end{tikzcd}
\]

Universal property of the product implies $\phi\psi=\id_{X\times Y}$, similarly universal property of fiber product implies $\psi\phi=\id_{X\times_Z Y}$.

\section{1 -- 1.2.Q -- Towers of Cartisian Diagrams is Cartisian}
% id: 1-1-2-q-towers-of-cartisian-diagrams-is-cartisia
% author: Tianyi
% updated: 2026-08-30

Here's the complete diagram:
% https://q.uiver.app/#q=WzAsNyxbMSwxLCJVIl0sWzIsMSwiViJdLFsxLDIsIlciXSxbMiwyLCJYIl0sWzEsMywiWSJdLFsyLDMsIloiXSxbMCwwLCJXJyJdLFswLDEsImYiXSxbMiwzLCJnIl0sWzEsMywiXFxnYW1tYSJdLFswLDIsIlxcYWxwaGEiXSxbMiw0LCJcXGJldGEiXSxbMyw1LCJcXGRlbHRhIl0sWzQsNSwiaCJdLFs2LDQsIlxcbXUiLDIseyJjdXJ2ZSI6M31dLFs2LDEsIlxcbnUiLDAseyJjdXJ2ZSI6LTJ9XSxbNiwyLCJcXHJobyIsMSx7InN0eWxlIjp7ImJvZHkiOnsibmFtZSI6ImRhc2hlZCJ9fX1dLFs2LDAsIlxcZXhpc3QhXFx0YXUiLDAseyJzdHlsZSI6eyJib2R5Ijp7Im5hbWUiOiJkYXNoZWQifX19XV0=
\[
\begin{tikzcd}[cramped]
	{W'} &  &  \\
	 & {U} & {V} \\
	 & {W} & {X} \\
	 & {Y} & {Z}
	\arrow["{f}", from=2-2, to=2-3]
	\arrow["{g}", from=3-2, to=3-3]
	\arrow["{\gamma}", from=2-3, to=3-3]
	\arrow["{\alpha}", from=2-2, to=3-2]
	\arrow["{\beta}", from=3-2, to=4-2]
	\arrow["{\delta}", from=3-3, to=4-3]
	\arrow["{h}", from=4-2, to=4-3]
	\arrow["{\mu}"', curve={height=18pt}, from=1-1, to=4-2]
	\arrow["{\nu}", curve={height=-12pt}, from=1-1, to=2-3]
	\arrow["{\rho}", dashed, from=1-1, to=3-2]
	\arrow["{\exist!\tau}", dashed, from=1-1, to=2-2]
\end{tikzcd}
\]

Let me explain: Suppose we are given maps $\mu,\nu$. Note we have a map $\gamma\nu: W'\to X$, this gives a unique map $\rho: W'\to W$ because the lower square is Cartesian. Now we have map $\rho: W'\to W, \nu: W'\to V$, so this gives a unique map $\tau: W'\to U$ because the upper square is also Cartesian.

To verify the tower is indeed Cartesian, we need to show $\nu=\tau f$ (checked because upper square is Cartesian) and $\mu=\beta\alpha\tau$. Chasing through the construction shows that $\beta\alpha\tau=\beta\rho=\mu$, as desired.

\section{1 -- 1.2.S -- Diagonal Base Change Diagram}
% id: 1-1-2-s-diagonal-base-change-diagram-incomplete
% author: Tianyi
% updated: 2026-08-30

This is an example of how Yoneda lemma makes abstract things concrete. Yoneda lemma implies:

$\textbf{Lemma}.$ A square in any category $\ms C$ is Cartesian iff applying $h_S=\Hom(S,-)$ of it gives a Cartesian square:
% https://q.uiver.app/#q=WzAsOCxbMCwwLCJBIl0sWzEsMCwiQiJdLFswLDEsIkMiXSxbMSwxLCJEIl0sWzMsMCwiaF9TKEEpIl0sWzQsMCwiaF9TKEIpIl0sWzMsMSwiaF9TKEMpIl0sWzQsMSwiaF9TKEQpIl0sWzAsMV0sWzAsMl0sWzIsM10sWzEsM10sWzQsNV0sWzUsN10sWzQsNl0sWzYsN11d
\[
\begin{tikzcd}[cramped]
	{A} & {B} &  & {h_S(A)} & {h_S(B)} \\
	{C} & {D} &  & {h_S(C)} & {h_S(D)}
	\arrow[from=1-1, to=1-2]
	\arrow[from=1-1, to=2-1]
	\arrow[from=2-1, to=2-2]
	\arrow[from=1-2, to=2-2]
	\arrow[from=1-4, to=1-5]
	\arrow[from=1-5, to=2-5]
	\arrow[from=1-4, to=2-4]
	\arrow[from=2-4, to=2-5]
\end{tikzcd}
\]

Next we show

$\textbf{Lemma}.$ $\Hom(S,A\times_C B)=\Hom(S,A)\times_{\Hom(S,C)}\Hom(S,B)$.

Proof: Note that $\Hom(S,A\times_C B)$ is the collection of data $(f:S\to A, g: S\to B)$ that agrees when maps to $C$ using $A\to C, B\to C$. But this is precisely the right hand side, where $Hom(S,A)\to \Hom(S,C)$ is by postcomposing $A\to C$, and similarly for $\Hom(S,B)\to \Hom(S,C)$. $\square$

Finally, applying the results above, our problem is equivalent to showing that the square below is Cartesian:
% https://q.uiver.app/#q=WzAsNCxbMCwwLCJoKFhfMSlcXHRpbWVzX3toKFkpfWgoWF8yKSJdLFsxLDAsImgoWF8xKVxcdGltZXNfe2goWil9aChYXzIpIl0sWzAsMSwiaChZKSJdLFsxLDEsImgoWSlcXHRpbWVzX3toKFopfWgoWSkiXSxbMCwxXSxbMSwzXSxbMCwyXSxbMiwzXV0=
\[
\begin{tikzcd}[cramped]
	{h(X_1)\times_{h(Y)}h(X_2)} & {h(X_1)\times_{h(Z)}h(X_2)} \\
	{h(Y)} & {h(Y)\times_{h(Z)}h(Y)}
	\arrow[from=1-1, to=1-2]
	\arrow[from=1-2, to=2-2]
	\arrow[from=1-1, to=2-1]
	\arrow[from=2-1, to=2-2]
\end{tikzcd}
\]

Up to relabeling, we just need to show the diagram
% https://q.uiver.app/#q=WzAsNCxbMCwwLCJYXzFcXHRpbWVzX1kgWF8yIl0sWzEsMCwiWF8xXFx0aW1lc19aIFhfMiJdLFswLDEsIlkiXSxbMSwxLCJZXFx0aW1lc19aIFkiXSxbMCwxXSxbMSwzXSxbMCwyXSxbMiwzXV0=
\[
\begin{tikzcd}[cramped]
	{X_1\times_Y X_2} & {X_1\times_Z X_2} \\
	{Y} & {Y\times_Z Y}
	\arrow[from=1-1, to=1-2]
	\arrow[from=1-2, to=2-2]
	\arrow[from=1-1, to=2-1]
	\arrow[from=2-1, to=2-2]
\end{tikzcd}
\]
is Cartesian but now everything are sets, and we know what fiber product in set means. This is a routine exercise to check.

\section{1 -- 1.2.T -- Coproduct is Sets is Disjoint Union}
% id: 1-1-2-t-coproduct-is-sets-is-disjoint-union
% author: Tianyi
% updated: 2026-08-30

Recall that in $Sets$, elements in $X_1\coprod X_2$ are given by elements of the form $(x_1,1),(x_2,2)$ where $x_1\in X_1, x_2\in X_2$. Thus of course we have a diagram
% https://q.uiver.app/#q=WzAsNCxbMCwwLCJXIl0sWzEsMSwiWF8xXFxjb3Byb2QgWF8yIl0sWzEsMiwiWF8xIl0sWzIsMSwiWF8yIl0sWzIsMV0sWzMsMV0sWzEsMCwiXFxleGlzdHMhXFxyaG8iLDEseyJzdHlsZSI6eyJib2R5Ijp7Im5hbWUiOiJkYXNoZWQifX19XSxbMiwwLCJcXHBoaSJdLFszLDAsIlxccHNpIiwyXV0=
\[
\begin{tikzcd}[cramped]
	{W} &  &  \\
	 & {X_1\coprod X_2} & {X_2} \\
	 & {X_1} &
	\arrow[from=3-2, to=2-2]
	\arrow[from=2-3, to=2-2]
	\arrow["{\exists!\rho}", dashed, from=2-2, to=1-1]
	\arrow["{\phi}", from=3-2, to=1-1]
	\arrow["{\psi}"', from=2-3, to=1-1]
\end{tikzcd}
\]
where $\rho(x_1,1)=\phi(x_1)$ and $\rho(x_2,2)=\psi(x_2).$

\section{1 -- 1.2.U -- Tensor Product is Fibered Coproduct}
% id: 1-1-2-u-tensor-product-is-fibered-coproduct
% author: Tianyi
% updated: 2026-08-30

We have the following diagram
% https://q.uiver.app/#q=WzAsNSxbMiwxLCJDIl0sWzEsMiwiQiJdLFsxLDEsIkJcXG90aW1lc19BIEMiXSxbMiwyLCJBIl0sWzAsMCwiUiJdLFszLDFdLFszLDBdLFswLDJdLFsxLDJdLFsxLDQsIlxcbXUiXSxbMCw0LCJcXG51IiwyXSxbMiw0LCJcXGV4aXN0IVxccmhvIiwxLHsic3R5bGUiOnsiYm9keSI6eyJuYW1lIjoiZGFzaGVkIn19fV1d
\[
\begin{tikzcd}[cramped]
	{R} &  &  \\
	 & {B\otimes_A C} & {C} \\
	 & {B} & {A}
	\arrow[from=3-3, to=3-2]
	\arrow[from=3-3, to=2-3]
	\arrow[from=2-3, to=2-2]
	\arrow[from=3-2, to=2-2]
	\arrow["{\mu}", from=3-2, to=1-1]
	\arrow["{\nu}"', from=2-3, to=1-1]
	\arrow["{\exist!\rho}", dashed, from=2-2, to=1-1]
\end{tikzcd}
\]
where the map $\rho: B\ot_A C\to R$ is defined by $\rho(b\ot c)=\mu(b)\ot \nu(c)$ and extended linearly.

\section{1 -- 1.3.C -- Limits in $Sets$}
% id: 1-1-3-c-limits-in-sets
% author: Tianyi
% updated: 2026-08-31

The set
\[
\left\{(a_i)_{i\in \ms I}\in\prod_{i\in \ms I}A_i: F(m)(a_j)=a_k\; \forall m\in \op{Mor}_{\ms I}(j,k)\subset \op{Mor}(\ms I)\right\}
\]
with the obvious projection $\pi_i$ to each $A_i$ is satisfies $\pi_k=F(m)\pi_j$ for every $m:j\to k$ by definition, and the universal property of the limit is satisfied by the above essentially due to the universal property of the product $\prod_{i\in \ms I}A_i$.

\section{1 -- 1.3.E -- Colimit in $Sets$}
% id: 1-1-3-e-colimit-in-sets
% author: Tianyi
% updated: 2026-09-01

The long equivalence relation in the quotient is to force the object to satisfy the coherence relation. The universal property of colimit follows from the universal property of coproduct.

\section{1 -- 1.4.C -- Tensor Hom Bijection}
% id: 1-1-4-c-tensor-hom-bijection
% author: Tianyi
% updated: 2026-09-01

We construct a bijection
\[
\Hom(M\ot N, P)\longleftrightarrow \Hom(M,\Hom(N,P))
\]
as follows: Given a map $M\ot N\to P$, composing with the natural map $M\times N\to M\ot N$ gives a bilinear map $M\times N\to P$, which is an element of $\Hom(M,\Hom(N,P))$.  Conversely, given an element in $\Hom(M,\Hom(N,P))$, i.e., a bilinear map $M\times N\to P$, universal property of tensor product gives a map $M\ot N\to P$. It is easy to verify that these two operations are inverses to each other.

\section{1 -- 1.4.D -- Tensor Hom Adjunction}
% id: 1-1-4-d-tensor-hom-adjunction
% author: Tianyi
% updated: 2026-09-01

We will only show one of the two commutative diagrams. Suppose we have a map $f:M\to M'$. We want to show the diagram
% https://q.uiver.app/#q=WzAsNCxbMCwwLCJcXG9wZXJhdG9ybmFtZXtIb219KE0nXFxvdGltZXMgTiwgUCkiXSxbMSwwLCJcXG9wZXJhdG9ybmFtZXtIb219KE1cXG90aW1lcyBOLFApIl0sWzAsMSwiXFxvcGVyYXRvcm5hbWV7SG9tfShNJyxcXG9wZXJhdG9ybmFtZXtIb219KE4sUCkpIl0sWzEsMSwiXFxvcGVyYXRvcm5hbWV7SG9tfShNLFxcb3BlcmF0b3JuYW1le0hvbX0oTixQKSkiXSxbMCwxXSxbMSwzXSxbMCwyXSxbMiwzXV0=
\[
\begin{tikzcd}[cramped]
	{\operatorname{Hom}(M'\otimes N, P)} & {\operatorname{Hom}(M\otimes N,P)} \\
	{\operatorname{Hom}(M',\operatorname{Hom}(N,P))} & {\operatorname{Hom}(M,\operatorname{Hom}(N,P))}
	\arrow[from=1-1, to=1-2]
	\arrow[from=1-2, to=2-2]
	\arrow[from=1-1, to=2-1]
	\arrow[from=2-1, to=2-2]
\end{tikzcd}
\]
commutes, where all the maps are the obvious ones. Let us chase through the definition in 1.4.C: Starting from the upper left corner suppose we have a map $\alpha: M'\ot N\to P$. Then going right gives a map $M\ot N\to M'\ot N\to P$ induced by $f:M\to M'$. Going down gives a bilinear map
\[M\times N\to M\ot N\to M'\ot N\to P.\]
This map is given by
\[
(m,n)\mapsto m\ot n\mapsto f(m)\ot n\mapsto \alpha(f(m)\ot n).
\]
Now if we first go down we get a map $M'\times N\to M'\ot N\to P$, then go right gives a map
\[
M\times N\to M'\times N\to M'\ot N\to P
\]
given by
\[
(m,n)\mapsto (f(m),n)\mapsto f(m)\ot n\mapsto \alpha(f(m)\ot n)
\]
which evidently agrees with the previous map displayed above. So the diagram commutes. The commutativity of the other diagram follows from a similar argument.

\section{1 -- 1.4.E -- $\Hom_A(N\ot_B A,M)\cong \Hom_B(N,M_B)$}
% id: 1-1-4-e-hom-a-n-ot-b-a-m-cong-hom-b-n-m-b
% author: Tianyi
% updated: 2026-09-01

We will only show that $\Hom_A(N\ot_B A,M)\cong \Hom_B(N,M_B)$ by constructing a bijection. I'll omit checking the naturality diagram...

Let $\phi: B\to A$ be the given change of coefficient map. Start with a map $\alpha: N\ot_B A\to M$. We construct a map $\beta: N\to M_B$ as follows: Let $\beta: n\mapsto n\ot_B 1_A\mapsto \alpha(n\ot_B 1_A).$

Conversely, given $\beta: N\to M_B$. We define $\alpha: N\ot_B A\to M$ as $\alpha: n\ot_B a\mapsto n\cdot\phi(a)\mapsto \beta(n\cdot \phi(a)).$ It is not hard to verify that these two operations are bijections.

\section{1 -- 1.5.A -- Two Exact Sequences of a Complex}
% id: 1-1-5-a-two-exact-sequences-of-a-complex
% author: Tianyi
% updated: 2026-09-01

The first exact sequence
\[
0\to \im f^{i}\to A^{i+1}\to \op{coker}f^{i}\to 0
\]
is simply the inclusion followed by quotient. That is, $\op{coker}f^{i}\cong A^{i+1}/\im f^{i}$ in an Abelian category. To see this, let $\alpha:A\to B$ be a morphism in an Abelian category. The definition we propose for cokernel satisfies the universal property diagram
% https://q.uiver.app/#q=WzAsNCxbMSwxLCJCIl0sWzIsMSwiQSJdLFswLDEsIkIvXFxvcGVyYXRvcm5hbWV7aW19XFxhbHBoYSJdLFswLDAsIlciXSxbMSwwLCJcXGFscGhhIiwyXSxbMCwyXSxbMCwzXSxbMSwzLCIwIiwyLHsiY3VydmUiOjF9XSxbMSwyLCIwIiwwLHsiY3VydmUiOi0yfV0sWzIsMywiXFxleGlzdCEiLDAseyJzdHlsZSI6eyJib2R5Ijp7Im5hbWUiOiJkYXNoZWQifX19XV0=
\[
\begin{tikzcd}[cramped]
	{W} &  &  \\
	{B/\operatorname{im}\alpha} & {B} & {A}
	\arrow["{\alpha}"', from=2-3, to=2-2]
	\arrow[from=2-2, to=2-1]
	\arrow[from=2-2, to=1-1]
	\arrow["{0}"', curve={height=6pt}, from=2-3, to=1-1]
	\arrow["{0}", curve={height=-12pt}, from=2-3, to=2-1]
	\arrow["{\exist!}", dashed, from=2-1, to=1-1]
\end{tikzcd}
\]
Hence is uniquely isomorphic to $\op{coker}\alpha$ by the universal property of the cokernel. The second exact sequence
\[
0\to H^i(A^\bullet)\to \op{coker}f^{i-1}\to \im f^i\to 0
\]
is inclusion $\ker f^i/\im f^{i-1}\to A_i/\im f^{i-1}$ composed with quotient since
\[
\frac{\op{coker}f^{i-1}}{H^i(A^\bullet)}=\frac{A_i/\im f^{i-1}}{\ker f^i/\im f^{i-1}}\cong \frac{A_i}{\ker f^i}\cong \im f^i
\]
by the third and first isomorphism theorems.

\section{1 -- 1.5.B -- Euler Characteristics of a Complex}
% id: 1-1-5-b-euler-characteristics-of-a-complex
% author: Tianyi
% updated: 2026-09-02

To show that $\sum_i (-1)^i\dim A^i=\sum_i (-1)^i h^i(A)$, we use the two exact sequences (1.5.6.3) in Vakil. This gives
\[\dim A^i=\dim \ker f^i+\dim \im f^i=\dim \im f^{i-1}+h^i(A)+\dim \im f^i.\]
Hence we obtain
\begin{align*}
\sum_i (-1)^i \dim A^i&=\sum_i (-1)^i h^i(A)+\sum_i (-1)^i [\dim \im f^i+\dim \im f^{i-1}]=\sum_i (-1)^i h^i(A).
\end{align*}
as the last sum cancels out due to a standard telescoping argument.

\section{1 -- 1.5.E -- Homotopic Maps Gives Same Map on Homology}
% id: 1-1-5-e-homotopic-maps-gives-same-map-on-homolog
% author: Tianyi
% updated: 2026-09-02

Since $f-g=dw+wd$ on the level of chains, and since every element $[\alpha]\in H^*(A)$ is represented by a closed (i.e. $d\alpha=0$) representative $\alpha\in A^*$, on the level of homology we have
\[
(f_*-g_*)([\alpha])=[dw\alpha+wd\alpha]=0,
\]
where the first term $dw\alpha=0$ is because it is a boundary, and the second term is zero because $d\alpha=0$ since $\alpha$ is closed.

\section{2 -- 2.1.A -- Germ is a Local Ring}
% id: 2-2-1-a-germ-is-a-local-ring
% author: Tianyi
% updated: 2026-09-03

We show $\mc O_p$ is local by showing $\mf m_p$ is the only maximal ideal. Let $(\varphi,U)$ be a representative of an element in $\mc O_p\setminus \mf m_p$, where $U$ is an open set containing $p$ on which $\varphi$ is nonzero (by continuity of $\varphi$). Then the image of $(\varphi^{-1},U)$ in $\mc O_p\setminus \mf m_p$ is an inverse of $(\varphi,U)$.

\section{2 -- 2.1.B -- The Cotangent Space from Germs Vanishing at a Point}
% id: 2-2-1-b-the-cotangent-space-from-germs-vanishing
% author: Anson
% updated: 2026-08-28

Consider
\[
\Phi:\mathfrak m_p\longrightarrow T_p^*X,
\qquad [f,U]\longmapsto (df)_p.
\]
It suffices to work in a coordinate chart, so we may assume that \(X=\mathbb R^n\) and \(p=0\).

To prove surjectivity, let \(\alpha\in T_0^*\mathbb R^n\) and define \(f_\alpha:\mathbb R^n\to\mathbb R\) by \(f_\alpha(x)=\alpha(x)\). Then \(f_\alpha\in\mathfrak m_0\) and \((df_\alpha)_0=\alpha\).

We now compute the kernel. If \(f\in\mathfrak m_0^2\), then by linearity it suffices to consider \(f=gh\) with \(g,h\in\mathfrak m_0\). The product rule gives
\[
(df)_0=h(0)(dg)_0+g(0)(dh)_0=0,
\]
so \(\mathfrak m_0^2\subseteq\ker\Phi\).

Conversely, suppose that \((df)_0=0\). After shrinking the domain of \(f\) around \(0\), for \(x\) in this neighborhood we have
\[
\begin{aligned}
f(x)
&=f(x)-f(0)\\
&=\int_0^1\frac{d}{dt}f(tx)\,dt\\
&=\sum_{i=1}^n x_i\int_0^1
\frac{\partial f}{\partial x_i}(tx)\,dt.
\end{aligned}
\]
Set
\[
g_i(x):=\int_0^1\frac{\partial f}{\partial x_i}(tx)\,dt.
\]
Then \(g_i(0)=\frac{\partial f}{\partial x_i}(0)=0\), so \(x_i,g_i\in\mathfrak m_0\). Since
\[
f=\sum_{i=1}^n x_i g_i,
\]
we have \(f\in\mathfrak m_0^2\). Therefore \(\ker\Phi=\mathfrak m_0^2\), and hence
\[
\mathfrak m_p/\mathfrak m_p^2\cong T_p^*X.
\]

\section{2 -- 2.2.A -- Presheaf is a Contravariant Functor}
% id: 2-2-2-a-presheaf-is-a-contravariant-functor
% author: Tianyi
% updated: 2026-09-12

Indeed, for every open set $U$ the functor assigns $\mc F(U)$, and for every inclusion $U\hookrightarrow{}V$ we have restriction map $\mc F(V)\to \mc F(U)$. The remaining properties are easy to verify.

\section{2 -- 2.2.D -- Function Spaces Form Sheaves}
% id: 2-2-2-d-function-spaces-form-sheaves
% author: Anson
% updated: 2026-08-28

The identity axiom is immediate in each case.

Let \(\{U_i\}\) be an open cover of \(U\), and let \(f_i\in\mathcal O(U_i)\) satisfy \(f_i|_{U_i\cap U_j}=f_j|_{U_i\cap U_j}\). They glue to a function \(f\) on \(U\): for \(x\in U\), choose \(i\) with \(x\in U_i\) and set \(f(x)=f_i(x)\). Compatibility makes \(f\) well-defined, and clearly \(f|_{U_i}=f_i\). Continuity, smoothness, and real-analyticity are local properties, so \(f\) has the required regularity. The same argument applies to continuous real-valued functions on an arbitrary topological space.

\section{2 -- 2.2.F -- Continuous Maps Form a Sheaf}
% id: 2-2-2-f-continuous-maps-form-a-sheaf
% author: Anson
% updated: 2026-08-28

Let \(\{U_i\}\) be an open cover of \(U\). For the identity axiom, if \(f,g:U\to Y\) satisfy \(f|_{U_i}=g|_{U_i}\) for all \(i\), then \(f=g\).

For gluability, let \(f_i:U_i\to Y\) be continuous maps agreeing on overlaps. Define \(f:U\to Y\) by \(f(x)=f_i(x)\) for \(x\in U_i\). This is well-defined, and \(f|_{U_i}=f_i\), so \(f\) is continuous. Thus continuous maps to \(Y\) form a sheaf.

\section{2 -- 2.2.H -- Pushforward Preserves Sheaves}
% id: 2-2-2-h-pushforward-preserves-sheaves
% author: Anson
% updated: 2026-08-28

Suppose that \(W\subseteq V\subseteq U\) are open in \(Y\). Since
\[
\pi^{-1}(W)\subseteq\pi^{-1}(V)\subseteq\pi^{-1}(U),
\]
the restriction maps of \(\mathcal F\) define those of \(\pi_*\mathcal F\), and the presheaf axioms follow from those of \(\mathcal F\).

If \(\{U_i\}\) is an open cover of \(U\), then \(\{\pi^{-1}(U_i)\}\) is an open cover of \(\pi^{-1}(U)\). Hence the identity and gluability axioms for \(\pi_*\mathcal F\) follow from those for \(\mathcal F\). Thus \(\pi_*\mathcal F\) is a sheaf whenever \(\mathcal F\) is.

\section{2 -- 2.2.I -- Pushforward Induces a Map on Stalks}
% id: 2-2-2-i-pushforward-induces-a-map-on-stalks
% author: Anson
% updated: 2026-08-28

For $[f,V]\in(\pi_*\mathcal F)_q$, where $q\in V\subseteq Y$ and $f\in\mathcal F(\pi^{-1}(V))$, the natural morphism is

\[
(\pi_*\mathcal F)_q\to \mathcal F_p,
\qquad
[f,V]\to [f,\pi^{-1}(V)].
\]

\section{2 -- 2.2.J -- The Stalk of a Module Sheaf}
% id: 2-2-2-j-the-stalk-of-a-module-sheaf
% author: Anson
% updated: 2026-08-28

For \([f,U]\in\mathcal O_{X,p}\) and \([a,V]\in\mathcal F_p\), define
\[
[f,U]\cdot[a,V]
=
\left[f|_{U\cap V}\,a|_{U\cap V},\,U\cap V\right].
\]
The module axioms follow from those on sections, so \(\mathcal F_p\) is an \(\mathcal O_{X,p}\)-module.

\section{2 -- 2.3.A -- Morphisms of (Pre)Sheaves Induce Morphisms of Stalks}
% id: 2-2-3-a-morphisms-of-pre-sheaves-induce-morphism
% author: Tianyi
% updated: 2026-09-12

Suppose we have a map $\phi:\mc F\to \mc G$. Let $[f,U]\in \mc F_p$ where $p\in U\subset X$, let $\phi_U: \mc F(U)\to \mc G(U)$. Then we define the induced map $\mc F_p\to \mc G_p$ is given as $[f,U]\mapsto [\phi_U(f),U]\in \mc G_p$. Note that this is well-defined, for if we use another representative we can always pass to common intersection of the defining open sets.

\section{2 -- 2.3.B -- Pushforward is Functorial}
% id: 2-2-3-b-pushforward-is-functorial
% author: Tianyi
% updated: 2026-09-12

Let $\pi: X\to Y$ be a continuous map. Suppose $\mc F$ is a sheaf on $X$, i.e., an object in $Sets_X$.

We first show that $\pi_*\mc F$ is a sheaf on $Y$. Gluability follows since if $\cup_{i\in I}U_i$ is an open cover of $Y$ then $\cup_{i\in I}\pi^{-1}(U_i)$ is an open cover of $X$ and we use gluability of $\mc F$ on $X$. Similarly identity axiom also holds.

Next, suppose we have morphism $\phi: \mc F\to \mc G$ of sheaves on $X$. Then we get $\pi_*\mc F\to \pi_*\mc G$ given by $\mc F(\pi^{-1}U)\to \mc G(\pi^{-1}U)$ for every open set $U\subset Y$. It is easy to check that this is indeed a morphism of sheaves.

\section{2 -- 2.3.C -- Sheaf  Hom}
% id: 2-2-3-c-sheaf-hom
% author: Tianyi
% updated: 2026-09-13

Let $\mc F,\mc G$ be two sheaves. We show that $\mc{Hom}(\mc F,\mc G)$ is a sheaf. By definition, for every open $U\hookrightarrow X$, the assignment is
\[
U\mapsto \op{Mor}(\mc F|_U,\mc G|_U).
\]
Given open containments $U\hookrightarrow V$, we have an evident restriction map $\op{Mor}(\mc F|_V,\mc G|_V)\to \op{Mor}(\mc F|_U,\mc G|_U)$ from the property of sheaf morphisms $\mc F|_V\to \mc G|_V$. These restrictions are compatible with each other, for if we have opens $U\hookrightarrow V\hookrightarrow W$, then for $O\subset W, O'\subset V, O''\subset U$ open, we have a commutative diagram
% https://q.uiver.app/#q=WzAsNixbMCwwLCJcXG1hdGhjYWx7Rn0oTykiXSxbMSwwLCJcXG1hdGhjYWwgRyhPKSJdLFswLDEsIlxcbWF0aGNhbHtGfShPJykiXSxbMCwyLCJcXG1hdGhjYWx7Rn0oTycnKSJdLFsxLDEsIlxcbWF0aGNhbCBHKE8nKSJdLFsxLDIsIlxcbWF0aGNhbCBHKE8nJykiXSxbMCwxLCJcXHBoaV9PIl0sWzIsNCwiXFxwaGlfe08nfSJdLFszLDUsIlxccGhpX3tPJyd9Il0sWzAsMiwiXFxyaG9fe08sTyd9IiwyXSxbMiwzLCJcXHJob197TycsTycnfSIsMl0sWzEsNCwiXFxyaG9fe08sTyd9Il0sWzQsNSwiXFxyaG9fe08nLE8nJ30iXSxbMCwzLCJcXHJob197TyxPJyd9IiwyLHsiY3VydmUiOjV9XSxbMSw1LCJcXHJob197TyxPJyd9IiwwLHsiY3VydmUiOi01fV1d
\[
\begin{tikzcd}[cramped]
	{\mathcal{F}(O)} & {\mathcal G(O)} \\
	{\mathcal{F}(O')} & {\mathcal G(O')} \\
	{\mathcal{F}(O'')} & {\mathcal G(O'')}
	\arrow["{\phi_O}", from=1-1, to=1-2]
	\arrow["{\phi_{O'}}", from=2-1, to=2-2]
	\arrow["{\phi_{O''}}", from=3-1, to=3-2]
	\arrow["{\rho_{O,O'}}"', from=1-1, to=2-1]
	\arrow["{\rho_{O',O''}}"', from=2-1, to=3-1]
	\arrow["{\rho_{O,O'}}", from=1-2, to=2-2]
	\arrow["{\rho_{O',O''}}", from=2-2, to=3-2]
	\arrow["{\rho_{O,O''}}"', curve={height=30pt}, from=1-1, to=3-1]
	\arrow["{\rho_{O,O''}}", curve={height=-30pt}, from=1-2, to=3-2]
\end{tikzcd}
\]
Here, $\rho$'s are restriction maps. Hence it is not hard to verify that the sheaf Hom is a presheaf.

Next we show that it is in fact a sheaf. This is where we use $\mc F,\mc G$ are sheaves: To prove gluability: Let $U\hookrightarrow X$ be open and $\cup_{i\in I}U_i=U$ is an open cover of $U$. Let $\phi_i\in \op{Mor}(\mc F|_{U_i},\mc G|_{U_i})$ be compatible morphisms. Then for any $O\hookrightarrow U$ open, and for any $\alpha\in \mc F(O)$, since $\phi_i|_{O\cap U_i\cap U_j}(\alpha)=\phi_j|_{O\cap U_i\cap U_j}(\alpha)\in \mc G(O\cap U_i\cap U_j)$ for all $i,j\in I$, these glue to a unique section $\phi_O(\alpha)\in \mc G(O)$ since $\mc G$ is a sheaf. This defines $\phi_O$ uniquely, proving gluability. Identity axiom can be proved similarly.

\section{2 -- 2.3.E -- Kernel Presheaf}
% id: 2-2-3-e-kernel-presheaf
% author: Tianyi
% updated: 2026-09-13

Let $\phi:\mc F\to \mc G$ be a morphism of presheaves. As suggested we define $\ker\phi: U\mapsto \ker\phi(U)$. To show that restriction maps are defined, we follow the hint and consider the below diagram:
% https://q.uiver.app/#q=WzAsOCxbMCwwLCIwIl0sWzEsMCwiXFxrZXIgXFxwaGkoVikiXSxbMiwwLCJcXG1hdGhjYWx7Rn0oVikiXSxbMywwLCJcXG1hdGhjYWwgRyhWKSJdLFswLDEsIjAiXSxbMSwxLCJcXGtlciBcXHBoaShVKSJdLFsyLDEsIlxcbWF0aGNhbCBGKFUpIl0sWzMsMSwiXFxtYXRoY2FsIEcoVSkiXSxbMCwxXSxbMSwyXSxbMiwzLCJcXHBoaShWKSJdLFs2LDcsIlxccGhpKFUpIiwyXSxbNSw2XSxbNCw1XSxbMiw2LCJcXHJob197VixVfSIsMl0sWzMsNywiXFxyaG9fe1YsVX0iXSxbMSw1LCJcXGV4aXN0ISIsMix7InN0eWxlIjp7ImJvZHkiOnsibmFtZSI6ImRhc2hlZCJ9fX1dXQ==
\[
\begin{tikzcd}[cramped]
	{0} & {\ker \phi(V)} & {\mathcal{F}(V)} & {\mathcal G(V)} \\
	{0} & {\ker \phi(U)} & {\mathcal F(U)} & {\mathcal G(U)}
	\arrow[from=1-1, to=1-2]
	\arrow[from=1-2, to=1-3]
	\arrow["{\phi(V)}", from=1-3, to=1-4]
	\arrow["{\phi(U)}"', from=2-3, to=2-4]
	\arrow[from=2-2, to=2-3]
	\arrow[from=2-1, to=2-2]
	\arrow["{\rho_{V,U}}"', from=1-3, to=2-3]
	\arrow["{\rho_{V,U}}", from=1-4, to=2-4]
	\arrow["{\exist!}"', dashed, from=1-2, to=2-2]
\end{tikzcd}
\]
The map is constructed as follows: For every $\alpha\in \ker\phi(V)$, view $\alpha\in \ker \phi(V)\hookrightarrow \mc F(V)$. Since the rightmost square commutes and the rows are exact, $\phi(U)\circ \rho_{V,U}(\alpha)$, hence $\rho_{V,U}(\alpha)$ is in the image of the map $\ker\phi(U)\to \mc F(U)$, i.e., there exists some $\beta\in \ker\phi(U)$ that maps to $\rho_{V,U}(\alpha)$. Then define the morphism to be $\alpha\mapsto \beta$. It is a routine verification that this map is unique and the diagram commutes.

To show the composition of restriction satisfies the desired property, let $U\hookrightarrow V\hookrightarrow W$ be inclusion of opens. We have the following diagram with commutative squares and exact rows:

% https://q.uiver.app/#q=WzAsMTIsWzAsMSwiMCJdLFsxLDEsIlxca2VyIFxccGhpKFYpIl0sWzIsMSwiXFxtYXRoY2Fse0Z9KFYpIl0sWzMsMSwiXFxtYXRoY2FsIEcoVikiXSxbMCwyLCIwIl0sWzEsMiwiXFxrZXIgXFxwaGkoVSkiXSxbMiwyLCJcXG1hdGhjYWwgRihVKSJdLFszLDIsIlxcbWF0aGNhbCBHKFUpIl0sWzAsMCwiMCJdLFsxLDAsIlxca2VyXFxwaGkoVykiXSxbMiwwLCJcXG1hdGhjYWwgRihXKSJdLFszLDAsIlxcbWF0aGNhbCBHKFcpIl0sWzAsMV0sWzEsMl0sWzUsNl0sWzQsNV0sWzIsNiwiXFxyaG9fe1YsVX0iLDJdLFszLDcsIlxccmhvX3tWLFV9Il0sWzgsOV0sWzksMTBdLFsxMCwxMV0sWzksMV0sWzEwLDIsIlxccmhvX3tXLFZ9IiwyXSxbMTEsMywiXFxyaG9fe1csVn0iXSxbMSw1XSxbMiwzXSxbNiw3XV0=
\[
\begin{tikzcd}[cramped]
	{0} & {\ker\phi(W)} & {\mathcal F(W)} & {\mathcal G(W)} \\
	{0} & {\ker \phi(V)} & {\mathcal{F}(V)} & {\mathcal G(V)} \\
	{0} & {\ker \phi(U)} & {\mathcal F(U)} & {\mathcal G(U)}
	\arrow[from=2-1, to=2-2]
	\arrow[from=2-2, to=2-3]
	\arrow[from=3-2, to=3-3]
	\arrow[from=3-1, to=3-2]
	\arrow["{\rho_{V,U}}"', from=2-3, to=3-3]
	\arrow["{\rho_{V,U}}", from=2-4, to=3-4]
	\arrow[from=1-1, to=1-2]
	\arrow[from=1-2, to=1-3]
	\arrow[from=1-3, to=1-4]
	\arrow[from=1-2, to=2-2]
	\arrow["{\rho_{W,V}}"', from=1-3, to=2-3]
	\arrow["{\rho_{W,V}}", from=1-4, to=2-4]
	\arrow[from=2-2, to=3-2]
	\arrow[from=2-3, to=2-4]
	\arrow[from=3-3, to=3-4]
\end{tikzcd}
\]

Now the last two columns composes as $\rho_{W,U}$, hence the composition of the restriction maps of kernels must be the restriction map of kernel $\ker\phi(W)\to \ker\phi(U)$ by uniqueness of restriction maps.

\section{2 -- 2.3.I -- Kernel Sheaf}
% id: 2-2-3-i-kernel-sheaf
% author: Tianyi
% updated: 2026-09-13

If $\phi:\mc F\to \mc G$ is a morphism of sheaves, then $\ker\phi$ is a sheaf: We will verify gluability since identity follows similarly. Let $U=\cup_{i\in I}U_i$. Let $\alpha_i\in \ker\phi(U_i)$ be compatible sections. Then using $\ker\phi(U_i)\hookrightarrow \mc F(U_i)$ we get compatible sections $\alpha_i\in \mc F(U_i)$. Hence they glue to a unique section $\alpha\in \mc F(U)$ since $\mc F$ is a sheaf. We show that this is indeed an element of the kernel: Since $\phi(\alpha)\in \mc G(U)$ satisfies that $\phi(\alpha)|_{U_i}=\phi|_{U_i}(\alpha|_{U_i})=0\in \mc G(U_i)$, we conclude that $\phi(\alpha)=0$ since $\mc G$ is a sheaf. Hence $\alpha\in \ker \phi(U)$. It is clear by construction that $\alpha|_{U_i}=\alpha_i$.

\section{2 -- 2.3.J -- Exponential Presheaf Sequence}
% id: 2-2-3-j-exponential-presheaf-sequence
% author: Tianyi
% updated: 2026-09-13

We first show that $0\to \underline{\Z}\to \mc O_X\to \mc F\to 0$ is an exact sequence of presheaves of Abelian groups. By definition, this means for every $U\subset X$ open the section of the sequence over $U$ is exact. By definition,
    \[
    \mc F(U):=\{f:U\to \C^*\text{ holomorphic}: \exists g: U\to \C\text{ holomorphic s.t.}\exp(2\pi i g)=f\}.
    \]
    Note that $\mc O_X(U)\to \mc F$ surjects because $g\mapsto f$ for any $f\in \mc F(U)$. Clearly $\underline{\Z}(U)\to \mc O_X(U)$ is an injection. Finally, we note that by standard results in complex analysis, a holomorphic function $g:U\to \C$ satisfies $\exp(2\pi i g)=0$ iff $g$ is integer valued, which shows exactness of the sequence at $\exp(2\pi i(-)):\mc O_X(U)\to \mc F$.

The presheaf $\mc F$ is not a sheaf because it does not satisfy gluability axiom: For example, I claim that the function $f(z)=z$ defined on $\C^*$ does not admit a global holomorphic log.

    For the sake of contradiction, let's assume $g$ is a globally defined holomorphic log of $f$. That is, $\exp(g(z))=z$. Differentiating gives $1=\exp(g(z))g'(z)=zg'(z)$, so $g'(z)=1/z$. But since $g$ is holomorphic, $\bar \del g=0$. Hence $dg=(\del+\bar\del)g=\del g=g'(z)dz$. Therefore around any closed loop $\gamma\subset \C^*$ we must have
    \[
    \int_\gamma g'(z)dz=\int_\gamma dg=\int_{\del \gamma}g=0
    \]
    by Stokes' theorem. However, if you take $\gamma$ to be the unit circle $C$ in $\C^*$ then
    \[
    \int_\gamma g'(z)dz=\int_C\frac{dz}{z}=\int_0^{2\pi}\frac{ie^{i\theta}d\theta}{e^{i\theta}}=2\pi i\neq 0.
    \]
    This is the desired contradiction.

\section{2 -- 2.4.A -- Sections are Determined by Germs}
% id: 2-2-4-a-sections-are-determined-by-germs
% author: Tianyi
% updated: 2026-09-13

We show that the natural map $\mc F(U)\to \prod_{p\in U}\mc F_p$ defined by $s\mapsto (s_p)_{p\in U}$ is injective. Suppose we have two sections $s,t\in \mc F(U)$ such that $(s_p)=(t_p)$. Then for every $p$, we can find open sets $U_p\hookrightarrow U$ such that $s|_{U_p}=t|_{U_p}$. But note that $\{U_p\}_{p\in U}$ for an open cover of $U$. Hence by the identity axiom, we conclude that $s=t$.

\section{2 -- 2.4.B -- Compatible Germs Come From Sections}
% id: 2-2-4-b-compatible-germs-come-from-sections
% author: Tianyi
% updated: 2026-09-13

Let $(s_p)_{p\in U}\in \prod_{p\in U}\mc F_p$ be a compatible germ. Let $U_p$ be the open containing $p$ in $U$ such that there is some $\tilde s_p\in \mc F(U_p)$ such that the germ of $\tilde s_p$ at all $q\in U_p$ is $s_q$. Now suppose we have $\td s_{p'}\in \mc F(U_{p'})$, with $U_p\cap U_{p'}\neq \emptyset$ and $p''\in U_p\cap U_{p'}$ is an arbitrary point. Now since $\td s_p,\td s_{p'}$ agrees in a neighborhood of $p''$ and $p''\in U_p\cap U_{p'}$ is arbitrary, by the identity axiom of the sheaf $\mc F|_{U_p\cap U_{p'}}$, we see that $\td s_p, \td s_{p'}$ agrees on $U_p\cap U_{p'}$. Hence by gluability of the sheaf $\mc F|_U$ with open cover $\{U_p\}_{p\in U}$, we see that $(s_p)_{p\in U}\prod_{p\in U}\mc F_p$ must come from some section $s\in \mc F(U)$.

\section{2 -- 2.4.C -- Morphisms are Determined by Stalks}
% id: 2-2-4-c-morphisms-are-determined-by-stalks
% author: Tianyi
% updated: 2026-09-13

Indeed, let $\phi_1,\phi_2: \mc F\to \mc G$ be morphisms of presheaves and $\mc G$ is a sheaf, and these two maps induces the same maps on each stalk. The following diagram commutes by how we constructed morphism on stalks:
% https://q.uiver.app/#q=WzAsNCxbMCwwLCJcXG1hdGhjYWwgRihVKSJdLFsxLDAsIlxcbWF0aGNhbCBHKFUpIl0sWzAsMSwiXFxwcm9kX3twXFxpbiBVfVxcbWF0aGNhbCBGX3AiXSxbMSwxLCJcXHByb2Rfe3BcXGluIFV9XFxtYXRoY2FsIEdfcCJdLFswLDFdLFswLDJdLFsxLDMsIiIsMCx7InN0eWxlIjp7InRhaWwiOnsibmFtZSI6Imhvb2siLCJzaWRlIjoiYm90dG9tIn19fV0sWzIsM11d
\[
\begin{tikzcd}[cramped]
	{\mathcal F(U)} & {\mathcal G(U)} \\
	{\prod_{p\in U}\mathcal F_p} & {\prod_{p\in U}\mathcal G_p}
	\arrow[from=1-1, to=1-2]
	\arrow[from=1-1, to=2-1]
	\arrow[hook, from=1-2, to=2-2]
	\arrow[from=2-1, to=2-2]
\end{tikzcd}
\]
Suppose $s\in \mc F(U)$. Then $(\phi_1(s)_p)=(\phi_2(s)_p)\in \prod_{p\in U}\mc G_p$, and we already showed the right vertical map is injective, so $\phi_1(s)=\phi_2(s)$, hence $\phi_1=\phi_2$ as desired.

\section{2 -- 2.4.D -- Isomorphisms are Determined by Stalks}
% id: 2-2-4-d-isomorphisms-are-determined-by-stalks
% author: Tianyi
% updated: 2026-09-13

Let $\phi:\mc F\to \mc G$ be a morphism of sheaves. We show that $\phi:\mc F\to \mc G$ is an isomorphism iff it induces isomorphisms on every stalk. First, if $\phi:\mc F\to \mc G$ is an isomorphism then clearly it induces isomorphism on each stalk.

Conversely, suppose $\phi:\mc F\to \mc G$ induces isomorphism on every stalk. Let $U\hookrightarrow X$ be any open set. Exercise 2.4.C already shows that $\phi(U):\mc F(U)\to \mc G(U)$ is injective, hence it remains to show that it is surjective. To this end, let $t\in \mc G(U)$ be a section with $(t_p)\in \prod_{p\in U}\mc G_p$ be the induced element on the stalks. We have a commutative diagram with vertical maps being injections:
% https://q.uiver.app/#q=WzAsNCxbMCwwLCJcXG1hdGhjYWwgRihVKSJdLFsxLDAsIlxcbWF0aGNhbCBHKFUpIl0sWzAsMSwiXFxwcm9kX3twXFxpbiBVfVxcbWF0aGNhbCBGX3AiXSxbMSwxLCJcXHByb2Rfe3BcXGluIFV9XFxtYXRoY2FsIEdfcCJdLFswLDFdLFswLDJdLFsyLDNdLFsxLDNdXQ==
\[
\begin{tikzcd}[cramped]
	{\mathcal F(U)} & {\mathcal G(U)} \\
	{\prod_{p\in U}\mathcal F_p} & {\prod_{p\in U}\mathcal G_p}
	\arrow[from=1-1, to=1-2]
	\arrow[from=1-1, to=2-1]
	\arrow[from=2-1, to=2-2]
	\arrow[from=1-2, to=2-2]
\end{tikzcd}
\]
Since $\prod_p \phi_p$ is an isomorphism, there exists $(s_p)\in \prod_p \mc F_p$ such that $\phi_p(s_p)=t_p$. What this means is that for every $p\in U$ there exists open sets $U_p\hookrightarrow U$, sections $\td s_p\in \mc F(U_p), \td t_p\in \mc G(U_p)$ such that one has $\phi(\td s_p)=\td t_p$.

We now claim that $(s_p)$ are compatible germs, hence coming from a section $s\in \mc F(U)$ by Exercise 2.4.B, and it follows that $\phi(s)=t$, proving surjectivity of $\phi$, and we will be done. To see compatibility, we claim that for every $q\in U_p$ we have $(\td s_p)_q=s_q$. We argue by contradiction: Assume $q\in U_p$ but $(\td s_p)_q\neq s_q$. But on the other hand $(\td s_q)_q=s_q$. Note that $U_p\cap U_q\neq \emptyset$ and $\td t_p=\td t_q=t$ on the intersection $U_p\cap U_q$. However we have $\phi_q((\td s_p)_q)\neq \phi_q((\td s_q)_q)$. This means that in any neighborhood of $q$ we have $\phi(\td s_p)\neq \phi(\td s_q)$, which means $\td t_p\neq \td t_q$ on any neighborhood of $q$, which is the desired contradiction.

\section{2 -- 2.4.F -- Sheafification is Unique}
% id: 2-2-4-f-sheafification-is-unique
% author: Tianyi
% updated: 2026-09-13

Suppose $\op{sh}':\mc F\to \mc F^{\op{sh'}}$ is another sheafification, then universal property gives two morphisms $\alpha,\beta$ between two sheafifications:
% https://q.uiver.app/#q=WzAsMyxbMCwwLCJcXG1hdGhjYWwgRiJdLFsxLDAsIlxcbWF0aGNhbCBGXntcXG9wZXJhdG9ybmFtZXtzaH19Il0sWzAsMSwiXFxtYXRoY2FsIEZee1xcb3BlcmF0b3JuYW1le3NoJ319Il0sWzAsMl0sWzAsMV0sWzEsMiwiXFxiZXRhIiwwLHsib2Zmc2V0IjotMSwic3R5bGUiOnsiYm9keSI6eyJuYW1lIjoiZGFzaGVkIn19fV0sWzIsMSwiXFxhbHBoYSIsMCx7Im9mZnNldCI6LTEsInN0eWxlIjp7ImJvZHkiOnsibmFtZSI6ImRhc2hlZCJ9fX1dXQ==
\[
\begin{tikzcd}[cramped]
	{\mathcal F} & {\mathcal F^{\operatorname{sh}}} \\
	{\mathcal F^{\operatorname{sh'}}} &
	\arrow[from=1-1, to=2-1]
	\arrow[from=1-1, to=1-2]
	\arrow["{\beta}", dashed, from=1-2, to=2-1]
	\arrow["{\alpha}", dashed, from=2-1, to=1-2]
\end{tikzcd}
\]
We claim $\alpha\beta=\id_{\mc F^{\op{sh}}}$. Indeed, commutativity of the diagram implies $\alpha\circ \op{sh'}=\op{sh}$. Hence we have the below diagram
% https://q.uiver.app/#q=WzAsNCxbMCwwLCJcXG1hdGhjYWwgRiJdLFsyLDAsIlxcbWF0aGNhbCBGXntcXG9wZXJhdG9ybmFtZXtzaH19Il0sWzEsMSwiXFxtYXRoY2FsIEZee1xcb3BlcmF0b3JuYW1le3NoJ319Il0sWzEsMiwiXFxtYXRoY2FsIEZee1xcb3BlcmF0b3JuYW1le3NofX0iXSxbMCwyLCJcXG9wZXJhdG9ybmFtZXtzaCd9IiwyXSxbMCwxLCJcXG9wZXJhdG9ybmFtZXtzaH0iXSxbMSwyLCJcXGJldGEiXSxbMiwzLCJcXGFscGhhIl0sWzAsMywiXFxvcGVyYXRvcm5hbWV7c2h9IiwyLHsiY3VydmUiOjJ9XSxbMSwzLCJcXG9wZXJhdG9ybmFtZXtpZH0iLDAseyJjdXJ2ZSI6LTJ9XV0=
\[
\begin{tikzcd}[cramped]
	{\mathcal F} &  & {\mathcal F^{\operatorname{sh}}} \\
	 & {\mathcal F^{\operatorname{sh'}}} &  \\
	 & {\mathcal F^{\operatorname{sh}}} &
	\arrow["{\operatorname{sh'}}"', from=1-1, to=2-2]
	\arrow["{\operatorname{sh}}", from=1-1, to=1-3]
	\arrow["{\beta}", from=1-3, to=2-2]
	\arrow["{\alpha}", from=2-2, to=3-2]
	\arrow["{\operatorname{sh}}"', curve={height=12pt}, from=1-1, to=3-2]
	\arrow["{\operatorname{id}}", curve={height=-12pt}, from=1-3, to=3-2]
\end{tikzcd}
\]
and the claim follows from uniqueness of the factorization. Similarly one can show $\beta\alpha=\id_{\mc F^{\op{sh'}}}$ as well. If $\mc F$ is a sheaf, then $\id: \mc F\to \mc F$ itself satisfies the said universal property of sheafification, hence $\mc F^{\op{sh}}=\mc F$ follows from uniqueness shown above.

\section{2 -- 2.4.G -- Sheafification is a Functor}
% id: 2-2-4-g-sheafification-is-a-functor
% author: Tianyi
% updated: 2026-09-13

Let $\phi: \mc F\to \mc G$ be a map of presheaves. Then composition with sheafification gives a map $\mc F\to \mc G\to \mc G^{\op{sh}}.$ Now by the universal property of sheafification $\mc F^{\op{sh}}$ we get a map $\phi^{\op{sh}}:\mc F^{\op {sh}}\to \mc G^{\op{sh}}$.
We note that this is functorial: If we have $\psi: \mc G\to \mc H$, then we have the following commutative diagram (in the category of presheaves):
% https://q.uiver.app/#q=WzAsNixbMCwwLCJcXG1hdGhjYWwgRiJdLFswLDEsIlxcbWF0aGNhbCBHIl0sWzAsMiwiXFxtYXRoY2FsIEgiXSxbMSwwLCJcXG1hdGhjYWwgRl57XFxvcGVyYXRvcm5hbWV7c2h9fSJdLFsxLDEsIlxcbWF0aGNhbCBHXntcXG9wZXJhdG9ybmFtZXtzaH19Il0sWzEsMiwiXFxtYXRoY2FsIEhee1xcb3BlcmF0b3JuYW1le3NofX0iXSxbMCwzXSxbMSw0XSxbMiw1XSxbMCwxLCJcXHBoaSJdLFsxLDIsIlxccHNpIl0sWzMsNCwiXFxwaGkgXntcXG9wZXJhdG9ybmFtZXtzaH19IiwyXSxbNCw1LCJcXHBzaSBee1xcb3BlcmF0b3JuYW1le3NofX0iLDJdLFswLDIsIlxccHNpXFxwaGkiLDIseyJjdXJ2ZSI6Mn1dLFszLDUsIlxccHNpIF57XFxvcGVyYXRvcm5hbWV7c2h9fSBcXHBoaV57XFxvcGVyYXRvcm5hbWV7c2h9fSIsMCx7ImN1cnZlIjotMn1dXQ==
\[
\begin{tikzcd}[cramped]
	{\mathcal F} & {\mathcal F^{\operatorname{sh}}} \\
	{\mathcal G} & {\mathcal G^{\operatorname{sh}}} \\
	{\mathcal H} & {\mathcal H^{\operatorname{sh}}}
	\arrow[from=1-1, to=1-2]
	\arrow[from=2-1, to=2-2]
	\arrow[from=3-1, to=3-2]
	\arrow["{\phi}", from=1-1, to=2-1]
	\arrow["{\psi}", from=2-1, to=3-1]
	\arrow["{\phi ^{\operatorname{sh}}}"', from=1-2, to=2-2]
	\arrow["{\psi ^{\operatorname{sh}}}"', from=2-2, to=3-2]
	\arrow["{\psi\phi}"', curve={height=12pt}, from=1-1, to=3-1]
	\arrow["{\psi ^{\operatorname{sh}} \phi^{\operatorname{sh}}}", curve={height=-12pt}, from=1-2, to=3-2]
\end{tikzcd}
\]
Hence we see by uniqueness that $(\psi\phi)^{\op{sh}}=\psi^{\op{sh}}\phi^{\op{sh}}$.

\section{2 -- 2.4.H -- Sheafification is a Sheaf}
% id: 2-2-4-h-sheafification-is-a-sheaf
% author: Tianyi
% updated: 2026-09-14

Given the definition of $\mc F^{\op{sh}}$ with the tautological restriction map, it is easy to see that $\mc F^{\op{sh}}$ is a presheaf. Now let $U\hookrightarrow X$ be an open set and $U=\cup_{i\in I}U_i$ an open cover. To see the identity axiom: If $(f_p)_{p\in U_i}=(g_p)_{p\in U_i}$ for every $i\in I$, of course $(f_p)_{p\in U}=(g_p)_{p\in U}$. To see the gluability axiom: Suppose we have $(f^i_p)_{p\in U_i}$ for every $i\in I$ such that $f^i_p=f^j_p$ for all $p\in U_i\cap U_j$, then clearly these glue to a global $(f_p)_{p\in U}$ that restricts to $(f^i_p)_{p\in U_i}$ on each $U_i$ and this is in fact a compatible germ by construction, since its restriction to each element of the open cover is.

\end{document}
